% V. Методика за експериментална проверка на еквивалентността.
% Ориентировъчен обем: 10–12 стр.
%
% Този раздел е САМОСТОЯТЕЛЕН ПРИНОС: измерва съвпадението на два независимо
% реализирани графични интерфейса — нещо, за което обичайно се приема, че не се
% измерва. Раздел VI ползва резултата му.
%
% Първоизточници в хранилището (числата се сверяват с тях преди предаване):
%   port/cpp/docs/comparison/tools/run_comparison.py           — разгръщане, задвижване, заснемане
%   port/cpp/docs/comparison/tools/vm_agent_{macos,windows}.py — агент в госта
%   port/cpp/tools/parity/pixel_score.py                       — SSIM + попикселна разлика, праговете
%   port/cpp/tools/parity/check_capture_integrity.py           — контрол на достоверността
%   port/CLAUDE.md, „Parity comparison policy“                 — правилата за оценяване
%   port/cpp/docs/comparison/PREDICTIONS.md                    — предварително регистрирани хипотези
\section{Методика за експериментална проверка на еквивалентността}
\label{sec:method}

Разделът превръща твърдението за независимост от платформения инструментариум в
измерима величина, от постановката на експеримента до мерките срещу недостоверно
измерване --- последните застъпени съзнателно най-обширно, защото основната заплаха
тук не е неточност, а правдоподобна неистина.

По вид подходът принадлежи към визуалното тестване на графичен
интерфейс~\cite{garousi2017vgt}, но с различна цел: там се сравняват две състояния на
една система във времето (регресия), тук --- две независими реализации на едно
описание в един и същ момент.

\subsection{Постановка на експеримента}
\label{subsec:exp_setup}

\textbf{Измерваната величина} е степента на съвпадение между изображението на
еталонната реализация и изображението на разработената, при идентично входно описание,
устройство, тема и момент.

\textbf{Хипотезата}: за интерфейс, описан изцяло в общото описание, съвпадението може
да се доведе до праг, отвъд който остатъчната разлика се обяснява само с изброими
причини --- отстраним дефект или изрично изключение (§\ref{subsec:rulings}).
Отхвърлянето ѝ би означавало неотстраним и необясним клас различия: граница на
независимостта извън обхвата на абстракцията.

\subsubsection{Едно описание, две независими реализации}

Валидността на сравнението почива на \term{един и същ файл}, не два еквивалентни: едни
и същи байтове захранват еталонното приложение и се вграждат при компилация в
разработената система --- разликата между изображенията не идва от входа, единствената
променлива е реализацията.

Измерването обхваща \boardPages{} интерфейса, подбрани да покрият всяка контрола и
всеки алгоритъм за разположение, не да представят типично приложение.

\textbf{Третата колона} не е излишна: разработената система има два входни пътя ---
програмен интерфейс и общо описание; невидяно разминаване между тях е дефект точно
колкото разминаването с еталона, затова тя изгражда същите интерфейси по другия път и
участва самостоятелно в оценяването.

\subsubsection{Обхват на измерването и различният статут на слоевете}

Измерването обхваща четири платформи, всяка с \boardPages{} интерфейса, две теми и две
колони --- четири оценки на интерфейс (§\ref{subsec:rulings}). Петата платформа, с
различния native инструментариум, участва само качествено: нейният слой изгражда
интерфейса от друго семейство контроли, отколкото еталонът на същата операционна
система, затова попикселното съвпадение там е принципно недостижимо. Изискването към
него е различно --- пълнота на съдържанието и съгласуваност между входните пътища, без
количествена таблица; оценка там, където няма смисъл, е по-лошо от липса на оценка.

\subsection{Автоматизирано заснемане}
\label{subsec:capture}

Конвейерът е показан на \figref{fig:capture_pipeline}: виртуална машина, управлявана
през отдалечена обвивка, чийто гост-агент разгръща, стартира, задвижва, представя
прозореца в точен размер и заснема; оценяването остава на управляващата машина, върху
изтеглените кадри.

Възпроизводимостта изисква средата да е записана: настолните платформи се снемат във
виртуална машина с изрично зададена разделителна способност 1512$\times$950 при мащаб
1:1, защото списъкът с режими съдържа и режими с двоен мащаб, при които кадърът излиза
двойно уголемен и отпада от проверката за размер, без това да личи от съобщенията ---
задаването на режима се отчита като успешно. Мобилните платформи се снемат на
симулатор и емулатор с фиксиран модел устройство.

Средите: домакин macOS 26.5.2 (управляваща машина и симулатор), виртуална машина с
macOS за настолните платформи, виртуална машина с Windows~11 за WinUI и Android
емулатор с API~34, срещу еталон закрепен на версия 10.0.71 --- версията, спрямо която
важи „предимство има изобразеното“ (§\ref{subsec:rulings}). \emph{Ограничение:}
версиите на гостовите операционни системи не са закрепени; при повторение те трябва да
бъдат отчетени наново, тъй като част от разминаванията (§\ref{subsec:taxonomy}, клас~5)
зависят от конкретната версия.

\begin{figure}[H]
\centering
\begin{tikzpicture}[font=\footnotesize, node distance=0.55cm,
  s/.style={draw, rounded corners=2pt, align=center, inner sep=4pt,
            text width=2.35cm, minimum height=0.8cm, fill=blue!5},
  d/.style={s, fill=green!10},
  q/.style={s, fill=orange!12},
  v/.style={s, fill=red!8},
  ar/.style={-Latex, thick}]

  \node[d] (xaml) {едно описание\\(общи байтове)};
  \node[s, below left=0.7cm and 0.9cm of xaml] (appA) {еталонно\\приложение};
  \node[s, below=0.7cm of xaml] (appB) {система:\\по код};
  \node[s, below right=0.7cm and 0.9cm of xaml] (appC) {система:\\от описанието};
  \draw[ar] (xaml) -- (appA);  \draw[ar] (xaml) -- (appB);  \draw[ar] (xaml) -- (appC);

  \node[s, below=0.8cm of appB] (launch) {ново стартиране\\за (интерфейс, тема)};
  \draw[ar] (appA) -- (launch); \draw[ar] (appB) -- (launch); \draw[ar] (appC) -- (launch);

  \node[s, right=0.9cm of launch] (shot) {представяне\\и заснемане};
  \node[q, right=0.9cm of shot] (chk) {проверка за\\достоверност};
  \node[s, below=0.7cm of chk] (score) {структурно сходство\\+ дял пиксели};
  \node[v, left=0.9cm of score] (verd) {4 оценки\\на интерфейс};

  \draw[ar] (launch) -- (shot);
  \draw[ar] (shot) -- (chk);
  \draw[ar] (chk) -- (score);
  \draw[ar] (score) -- (verd);
  \node[note, below=0.5cm of verd, xshift=1.6cm, text width=7.2cm, font=\scriptsize]
    {Кадър, който не премине проверката, се \textbf{отхвърля} и се заснема наново ---
     никога не се замества с предходния (§\ref{subsec:integrity}).};
\end{tikzpicture}
\caption{Гръбнакът на измерването. Едно описание захранва три приложения; всяка двойка
„интерфейс\,--\,тема“ се снема от новостартиран процес; кадър, който не премине проверката
за достоверност, се отхвърля и заснемането се повтаря, вместо да бъде заместено.}
\label{fig:capture_pipeline}
\end{figure}

\subsubsection{Изолация на всяко измерване}

Всяка двойка „интерфейс\,--\,тема“ се снема от \term{новостартиран процес}: интерфейсът
се задава чрез променлива на средата при стартиране, не чрез навигация в работещо
приложение --- решение заради заснемането, със следствие за Раздел~\ref{sec:results}:
всяка извадка идва от студен процес, независимо от реда на минаване.

Обратната страна е ограничение на инструмента: щом всеки интерфейс се наблюдава в
отделен процес, конвейерът \emph{не може} да наблюдава освобождаване на памет при
напускането му --- измерва следата, не своевременността на освобождаването ѝ.

\subsubsection{Задвижване на интерфейса}

Задействието има два пътя. По \emph{абсолютни координати} --- натискане в точка от
госта --- работи винаги, но е чувствително: контрола, изместена с няколко пиксела,
превръща натискането в удар по празно място, безшумно. По \emph{идентификатор на
елемент} --- вграден агент по локален мрежов канал намира елемента и предизвиква
действието --- устойчив е на разместване, но изисква идентификатори в описанието, а не
всеки интерфейс ги има.

Изборът е автоматичен: с идентификатор и отговор от агента се ползва вторият път,
иначе координати; отказът се отбелязва в дневника, за да не се скрие тиха замяна на
устойчив сценарий в чувствителен.

Всеки задвижван интерфейс носи собствен \term{сценарий} --- файл в
\code{docs/comparison/scenarios/}, изброяващ поредица от именувани стъпки (,,initial“ е
по договор кадърът в покой), всяка с незадължително действие (натискане, превъртане,
влачене, плъзгане или въвеждане на текст) и координати като дял от заснетата
повърхност. От интерфейсите на борда \textbf{38} носят такъв сценарий, повечето с
повече от една стъпка; останалите се заснемат само в покой.

Изчакването преди заснемане след действие не е адаптивно измерване на покой, а
фиксирана пауза, заменима поотделно за всяка стъпка. Стойността се калибрира ръчно,
чрез повторение: четири последователни стартирания на еталона върху \code{box\_view}/iOS
показват при пауза 0,6~s разсейване от 52~px между собствените кацания на еталона, а
при 2,0~s --- и осемте стартирания на двете страни кацат в един и същ ред пиксели
(коментар в \code{scenarios/box\_view.toml}). Паузата защитава срещу същия смущаващ
фактор като ,,кадър по време на преход“ (§\ref{subsec:integrity}) --- разликата е, че
тук се предотвратява, а не се открива post factum.

\subsection{Правила за оценяване}
\label{subsec:rulings}

Сравнението на две изображения дава число; превръщането му в „дефект“ или „допустимо
разминаване“ е решение --- вземано наново всеки път, то би позволило всяко неудобно
разминаване да се обясни post factum. Затова методиката включва \term{предварително
записан набор от правила}, всяко формулирано веднъж, датирано и приложено механично;
наборът расте само в една посока --- ново наблюдение се записва като въпрос и се решава
\emph{преди} да бъде приложено, не след резултата.

\subsubsection{Основното правило}

Еталонът е меродавен за цялото съдържание --- всяка разлика в цвят, размер, разстояние
или текст е дефект на разработената система, никога „несъвършенство на еталона“; без
това правило всяко разминаване би подлежало на предоговаряне и мярката би загубила
смисъл.

\subsubsection{Изрично изброените изключения}

Изключенията от правилото са именно изброени, не изведени по преценка:

\begin{enumerate}[topsep=2pt,itemsep=1pt,leftmargin=*]
\item \textbf{Рамка на заснемане.} Еталонното приложение поставя всеки интерфейс във
      вътрешна рамка с голям външен отстъп и подрязва горния и долния му край ---
      равномерно отместване на приложението-носител, не на изобразяването, което не
      изменя размер, цвят или вътрешен отстъп на нито една контрола; само отместването
      отпада от сметката, разлика в самите контроли се брои.
\item \textbf{Особеност на конкретен native инструментариум, която еталонът също
      търпи.} На една от настолните среди еталонът не изобразява точките на
      индикатор, определени вградени изображения и сянка, докато на другите две ги
      изобразява; системата ги изобразява навсякъде --- възпроизвеждането на липсата
      би било грешка, системата е по-вярна на описанието.
\item \textbf{Артефакт на измерването}, а не на изобразяването (§\ref{subsec:integrity}).
\end{enumerate}

Всяко изключение е \emph{ограничено} до конкретна контрола и среда, не общо разрешение
за разминаване --- изключение, което важи навсякъде, е отказ от мярката.

\subsubsection{Противоречие между документацията на еталона и поведението му}

Отделно правило урежда случая, в който изходният текст на еталона казва едно, а
действително изобразеното --- друго, защото изходният текст е снимка от момент, а
измерването е срещу разпространената версия. Правилото: \textbf{предимство има
изобразеното} --- изходният текст остава източник за \emph{механизма}, не за
стойност, променена оттогава; разминаването се отбелязва в кода с посочване на двата
източника.

\subsubsection{Четирите задължителни сравнения}

Еталонът се съпоставя \emph{независимо} с всяка от двете колони на разработената
система, в двете теми --- четири оценки на интерфейс:

\begin{center}
\begin{tabular}{@{}llll@{}}
\textbf{№} & \textbf{Еталон} & \textbf{Съпоставена колона} & \textbf{Тема} \\
\hline
1 & светла & изграждане по програмен интерфейс & светла \\
2 & светла & изграждане от описанието & светла \\
3 & тъмна  & изграждане по програмен интерфейс & тъмна \\
4 & тъмна  & изграждане от описанието & тъмна \\
\end{tabular}
\end{center}

Четирите оценки \textbf{не се осредняват}: това би скрило точно интересния случай ---
единият входен път верен, другият не; средното на „вярно“ и „невярно“ изглежда „почти
вярно“, погрешно описание на положението.

\subsection{Количествено оценяване}
\label{subsec:scoring}

Оценяването е детерминистично: две допълващи се мерки, изчислени за всяка двойка
изображения.

\textbf{Структурно сходство}~\cite{wang2004ssim} --- чувствително към промени в
структурата, слабо чувствително към равномерни отмествания на яркост; подходящо, защото
интересът е към разлики във \emph{формата и разположението}.

За два съответни прозореца $x$ и $y$ мярката е

\begin{equation}
\mathrm{SSIM}(x,y)=\frac{(2\mu_x\mu_y+C_1)\,(2\sigma_{xy}+C_2)}
                        {(\mu_x^{2}+\mu_y^{2}+C_1)\,(\sigma_x^{2}+\sigma_y^{2}+C_2)},
\label{eq:ssim}
\end{equation}

където $\mu$ е средната яркост, $\sigma^{2}$ --- дисперсията, $\sigma_{xy}$ ---
ковариацията между прозорците, а $C_1=(K_1L)^2$, $C_2=(K_2L)^2$ са стабилизиращи
константи при $L=255$, $K_1=0{,}01$, $K_2=0{,}03$; трите множителя отговарят на
\emph{яркост}, \emph{контраст} и \emph{структура}, а оценката за изображението е
средното на $\mathrm{SSIM}(x,y)$ по всички прозорци, върху яркостния канал. Равномерно
отместване на яркостта изменя $\mu$, но оставя $\sigma_{xy}$ и $\sigma$ непроменени,
така че структурният множител не се влияе --- изместване на контрола обаче руши
ковариацията и мярката пада рязко.

\textbf{Дял на видимо различаващите се пиксели.} Втората мярка е

\begin{equation}
D=\frac{1}{N}\sum_{i=1}^{N}
  \Big[\max_{c\in\{R,G,B\}}\big|a_i^{c}-b_i^{c}\big| > \tau\Big],
\label{eq:diff}
\end{equation}

където $N$ е броят пиксели, $a_i$, $b_i$ --- съответните пиксели, $[\cdot]$ е индикатор,
а $\tau$ е прагът, отделящ разлика от шум при заглаждане на шрифта; максимумът по
канали, не средното, е съзнателен --- разлика само в синьото например е разлика, а
средното я разрежда до незабележимост.

Двете се отчитат заедно: сходството остава високо при дребна, систематична разлика в
цвят, която делът на пикселите улавя; делът остава нисък при отместен, но значим
елемент, което сходството наказва --- всяка мярка поотделно има сляпо петно, заедно
нямат общо.

Праговете, приложени към \emph{по-лошата} от двете теми, не средната:

\begin{center}
\begin{tabular}{@{}lll@{}}
\textbf{Оценка} & \textbf{Структурно сходство} & \textbf{Различаващи се пиксели} \\
\hline
съвпада             & $\geq 0{,}98$ & и $\leq 1{,}0\,\%$ \\
незначителна разлика & $\geq 0{,}90$ & и $\leq 8{,}0\,\%$ \\
съществена разлика   & \multicolumn{2}{l}{всичко останало} \\
\end{tabular}
\end{center}

Двете условия важат \emph{едновременно}. Пикселът се брои различаващ се при разлика по
някой канал над \textbf{25} от 255 --- праг, поглъщащ и шума от заглаждане на шрифта,
тоест наказващ еднакво изобразен текст. Калибрирани са срещу две наблюдения: долната
граница за „съвпада“ държи заглаждането на шрифта под себе си при иначе идентично
разположение; горната за „незначителна“ пуска изместена или липсваща контрола над себе
си, като съществено разминаване.

Втора проверка --- по колко интерфейса всяка мярка поотделно би дала различна оценка
--- идва от \code{comparison.json}: от \textbf{688} двойки, при \textbf{8} самостоятелно
различна е оценката на пикселния дял, при \textbf{3} --- на сходството, никога и
двете едновременно. Малките, ненулеви числа \emph{в двете посоки} показват, че
премахването на която и да
е мярка би сменило оценката на единадесет интерфейса --- при разминаване само в
едната посока по-простата мярка щеше да стигне.

\subsubsection{Оценяване на подвижни интерфейси}
\label{subsec:motion_scoring}

Задвижена или анимирана страница не дава една двойка изображения, а две
\emph{последователности} от кадри --- по една на колона --- и мерките по-горе не казват
какво да се прави с втора, трета, четвърта двойка. Оценяването минава през пет стъпки:

\begin{enumerate}[topsep=2pt,itemsep=1pt,leftmargin=*]
\item \textbf{Сдвояване по стъпка.} Кадрите от двете колони се сдвояват по името на
      стъпката, която ги е произвела; кадър без чифт \emph{отпада}, вместо да се
      сдвои с най-близкия по ред --- отпадналите се преброяват и се отчитат изрично,
      вместо мълчаливо да намалят извадката.
\item \textbf{Времево подравняване.} Едно общо отместване (до 3 кадъра) компенсира
      разлика в скоростта на заснемане между двете приложения; отхвърля се, ако би
      задържало под 60\,\% от сдвоените кадри --- страница, обходена в грешен ред, не
      може да бъде подравнена към съгласие.
\item \textbf{Пространствено подравняване.} Отделно, ограничено до 48~px по вертикала,
      компенсира инерционно разминаване в кацането на превъртане --- прилага се
      \emph{само} ако кадърът в покой преди действието вече съвпада с борда, и никога
      не връща по-лоша оценка от неподравнената; изместване извън тази граница остава
      непоправено и се отчита като разлика.
\item \textbf{Собствено движение.} Всяка колона се сравнява \emph{сама със себе си}
      --- първият ѝ кадър срещу най-различния измежду останалите --- за да раздели
      ,,и двете страни стоят неподвижни“ от ,,еталонът анимира, системата замръзва“:
      кръстосаното структурно сходство не забелязва малък въртящ се индикатор върху
      иначе еднакъв кадър.
\item \textbf{Устойчив край, при неуспех на кръстосаното сравнение.} На платформата,
      на която движението \emph{не е} възпроизводимо между стартирания\footnote{Android
      --- измерено разсейване между собствени стартирания $\approx$0,65\,\%, срещу
      $\approx$53{,}3\,\% разлика между кръстосано сдвоени кадри на един и същ момент от
      жеста: 80 пъти по-голямо.}, сдвояването по стъпка е твърде шумно, за да реши
      само. Преди да съществуваше тази стъпка, такъв случай оставаше с горна граница
      жълто, никога зелено. Оценката пита вместо това дали \term{последните три кадъра}
      на всяка колона са спрели да се менят помежду си --- и само ако и двете страни са
      спрели, сравнява установеното им състояние едно с друго, с прощаването от точка~3.
      Страница, която никога не спира (например пулсиращ индикатор), остава без
      принудена оценка, вместо погрешно зелена или червена.
\end{enumerate}

Точка~5 заменя единствено \emph{оценката на движението} в записа на клетката --- не
цвета ѝ: класификацията на цвета никога не се повишава само от заключение за движение.
Разграничението пази основното правило на §\ref{subsec:rulings} --- еталонът остава
меродавен --- дори когато механизмът, потвърждаващ \emph{защо} дадено разминаване не е
дефект, е по-сложен от статично сравнение на два кадъра.

Класификацията разчита и на \term{еталонът срещу себе си}: едно и също, непроменено
приложение на еталона, стартирано неколкократно и задвижено по един и същ начин;
собственото му разсейване между стартирания служи за шумов праг --- разлика, не
по-голяма от него, не може да бъде приписана на портираната система. Сравненията се
групират по един и същ ангажимент от двете страни: по-ранен опит без това групиране
обърка прекомпилиране на системата между заснемания с истинска нестабилност и докладва
грешна цифра --- същият клас смущаващ фактор, който §\ref{subsec:integrity}
каталогизира за статичните кадри.

\subsubsection{Отказ от оценяване чрез езиков модел}

Първоначално оценките включваха и оценка чрез езиков модел --- \emph{изоставена}, след
като две независими такива оценки на едни и същи двойки се разминаха помежду си и с
измеримото. Причината е методологична, не техническа: невъзпроизводима оценка не служи
за мярка, колкото и полезно описание да дава. Отзивът остава полезен при
\emph{диагностика} на защо две изображения се различават, не при \emph{оценяване}.

\subsection{Достоверност на измерването}
\label{subsec:integrity}

Този параграф разглежда основната заплаха пред изследването: заснемащ конвейер, който
докладва „заснети 1104, откази 0“, твърди само, че файловете са записани, не че
показват онова, което претендират. Описаните по-долу режими на отказ са настъпвали в
хода на работата и нито един не е бил сигнализиран от самия конвейер.

\subsubsection{Заснемане, което произвежда правдоподобна неистина}

\textbf{Кадър от чужда колона.} Когато представянето на прозореца се провали, а
предходният файл остане на мястото си, изтеглянето връща предишния файл под новото
име --- валидно изображение, което сравнява реализацията със самата себе си или с
чужда тема. Установеният случай доведе до около 80\% съвпадение за интерфейс без
никакъв дефект, непоправимо с промяна в реализацията.

\textbf{Съвпадащи байтове между теми.} Ако темата не е приложена при стартиране,
„тъмното“ заснемане е второ светло --- признак, точен защото фонът винаги се различава
между теми: два байт-идентични кадъра на един интерфейс са сигнал за дефект, не
рядкост.

\textbf{Кадър по време на преход.} Заснемане преди установяване на изобразяването ---
и трите минават всяка проверка, основана само на наличие и размер на файла.

\subsubsection{Механични проверки вместо доверие}

Контролът е механичен и се извършва \emph{преди} оценяването:

\begin{itemize}[topsep=2pt,itemsep=1pt,leftmargin=*]
\item хеширане на всички кадри и откриване на съвпадащи байтове между колони (винаги
      дефект) и между теми в една колона (дефект, освен ако интерфейсът не е обявен
      изрично за независим от темата);
\item изискване неуспешното представяне да \emph{отпадне} като кадър, вместо да бъде
      заместено с предходния;
\item изискване двете колони на едно сравнение да произхождат от едно и също
      изпълнение.
\end{itemize}

Изведеното от това правило е формулирано отрицателно: \term{при съмнение кадърът се
отхвърля, а не се приема}. Всеки резервен вариант при заснемане произвежда измислица
--- разликата е само в това колко правдоподобна.

\subsubsection{Смущаващи фактори, установени в хода на работата}

Трите фактора по-долу първоначално са били диагностицирани като дефект на
реализацията, преди да бъдат разпознати.

\begin{enumerate}[topsep=2pt,itemsep=1pt,leftmargin=*]
\item \textbf{Смяна на денонощието между двете заснемания.} Интерфейс с дата или час
      изменя ширината на изобразения низ, имитирайки дефект в разположението ---
      съпоставими са само интерфейси, чиито колони идват от едно изпълнение.
\item \textbf{Различен мащаб на екрана.} Заснемане на еталона при увеличен системен
      мащаб поражда цял клас „еталонът е по-едър“, който не е дефект --- опит да се
      „поправи“ води до влошаване на реализацията.
\item \textbf{Източник на темата.} Темата произхожда от операционната система, не от
      настройка на приложението; без изрично управление двете колони се снемат в
      различни теми.
\end{enumerate}

Към същия списък спада системната лента на една от платформите: час и заряд се
различават между заснемания без общо с изобразяването --- изключва се чрез отрязване
на горната ивица преди оценяването, около 135 реда, определени чрез измерване, докато
съдържанието под нея се подравнява.

\subsection{Предварителна регистрация на хипотезите}
\label{subsec:preregistration}

Хипотезите за цената на реализацията са записани \emph{преди} първото измерване, заедно
с изискванията за валидност. Практиката е установена в емпиричната софтуерна
инженерия като \term{регистриран доклад}~\cite{ernst2023registered}; тук е приложена в
опростен вид --- регистрацията е част от документацията на хранилището, датирана.

Изискванията за валидност, договорени предварително:

\begin{itemize}[topsep=2pt,itemsep=1pt,leftmargin=*]
\item еднаква конфигурация на компилация от двете страни --- отстраняване на грешки
      срещу издаване е негодно сравнение в двете посоки;
\item повторения и разпределения, не единични числа;
\item времето за кадър по горни процентили, не по средна стойност --- забавянето е
      явление на опашката на разпределението;
\item разграничаване на студено от топло стартиране;
\item определение за „стартиране“, независимо от реализацията --- до първия заснет
      кадър, не началния екран, а не до появата на прозорец, различно предхождащ
      изобразяването в двете реализации;
\item неизмерим слой се отчита като \emph{неизмерен}, не се осреднява с останалите.
\end{itemize}

Регистрацията включва и хипотеза \textbf{срещу} очакването на автора --- за величина,
при която се очаква да \emph{не} се установи разлика; потвърждаването ѝ е резултат, не
пропуск --- изследване с изцяло сбъднати прогнози се чете като застъпничество.

Пълният текст, с датите и формулировките в оригиналния им вид, е приложен към работата
(Приложение~Д), възпроизведен без редакция --- преформулирана след факта прогноза не е
прогноза.
